\section{Phase III: Fabric Network Deployment and Integration Validation}
\label{sec:phase3_fabric_network}

The final implementation boundary deploys the unchanged \texttt{m4evidence} contract from
Section~\ref{sec:phase3_chaincode_unit} and submits the canonical \texttt{urllib3\_W03} evidence from
Section~\ref{sec:phase3_offchain_verifier} through an actual Fabric Gateway client. This step tests the properties
that local contract mocks cannot establish: organizational approval and endorsement, ordering and commit status,
authenticated identity handling, peer-visible state, channel history, and Gateway event delivery.

\subsection{Frozen environment and deployment}
\label{sec:phase3_fabric_deployment}

The run used Fabric \texttt{v2.5.15} (commit \texttt{83c7930}) from the Fabric 2.5 LTS line, Fabric CA binaries
\texttt{v1.5.17}, Docker Engine \texttt{29.4.3}, and Docker Compose \texttt{5.1.3}. The official
\texttt{fabric-samples} test network was detached at commit
\path{65592350d7d7c51b02c8a4d89383d4bbdcc45725}. Because the installer did not expose a matching
\texttt{fabric-samples v2.5.15} tag, recording this commit avoids treating a moving sample branch as a frozen input.
Fabric identities were generated locally with \texttt{cryptogen}; Fabric CA was installed as part of the frozen tool set
but no CA service was used. This wording distinguishes configured X.509 MSP identities from CA enrollment.

The topology contains \texttt{peer0.org1.example.com}, \texttt{peer0.org2.example.com}, and one Raft orderer. Both
peers joined \texttt{mychannel}, using GoLevelDB. Four byte-checked runtime files were staged from the locally tested
source and packaged for deployment. Table~\ref{tab:phase3_fabric_definition} records the committed definition.

\begin{table}[H]
\centering
\small
\renewcommand{\arraystretch}{1.08}
\begin{tabularx}{\textwidth}{p{4.2cm}X}
\toprule
\textbf{Deployment property} & \textbf{Measured value} \\
\midrule
Chaincode name / version / sequence & \texttt{m4evidence} / \texttt{1.0} / \texttt{1} \\
Package size / source files & 13,275 bytes / 4 \\
Package SHA-256 & \path{5e767ab94efbb4a63b8958f44249b0c0c2d7cede71738593c813282f091fabb7} \\
Organization approvals & \texttt{Org1MSP=true}; \texttt{Org2MSP=true} \\
Endorsement policy & \texttt{AND('Org1MSP.peer','Org2MSP.peer')} \\
Gateway / gRPC versions & \texttt{1.10.1} / \texttt{1.14.4} \\
Production dependency findings & 0 \\
\bottomrule
\end{tabularx}
\caption{Committed chaincode definition and client environment.}
\label{tab:phase3_fabric_definition}
\end{table}

The validator queried the committed definition on both peers and decoded its binary validation parameter. It requires a
threshold of two with the \texttt{Org1MSP PEER} and \texttt{Org2MSP PEER} principals. This check is stronger than
matching only a displayed policy string. The package was installed on both peers with identifier
\path{m4evidence_1.0:5e767ab94efbb4a63b8958f44249b0c0c2d7cede71738593c813282f091fabb7}.

\subsection{Authenticated positive and negative paths}
\label{sec:phase3_fabric_gateway_result}

The Gateway client used the generated \texttt{User1} certificate and private key for each organization. The accepted
proposal used the Org1 identity permitted by chaincode and explicitly requested endorsements from both organizations. The
17 integration checks all passed; Table~\ref{tab:phase3_fabric_gateway_checks} groups them by purpose.

\begin{table}[H]
\centering
\small
\renewcommand{\arraystretch}{1.08}
\begin{tabularx}{\textwidth}{p{3.6cm}Xr}
\toprule
\textbf{Group} & \textbf{Checks} & \textbf{Passed} \\
\midrule
Precondition and rejection & independent record ID, absent initial state, malformed JSON, one-ppm mutation, unauthorized Org2 submitter, and no state from rejected proposals & 6/6 \\
Commit and event & valid Org1 commit, exact 26-field proposed record, Gateway event reception, and exact event payload & 4/4 \\
Ledger observation & state existence, identical reads from both peers, project query, developer query, and exact history & 5/5 \\
Immutability & duplicate rejection and preservation of original state & 2/2 \\
\midrule
Total & & 17/17 \\
\bottomrule
\end{tabularx}
\caption{Authenticated Fabric Gateway integration checks.}
\label{tab:phase3_fabric_gateway_checks}
\end{table}

Malformed JSON, the one-ppm mutation, and the Org2 submission failed during endorsement and created no state. After the
valid commit, the project and developer indexes returned the same record, both peers returned byte-equivalent parsed state,
and history contained the single expected update. A second valid submission was rejected and did not alter that state.

\subsection{Committed transaction and independent validation}
\label{sec:phase3_fabric_transaction}

The accepted record retained the deterministic identifier already derived in the local tests:

\begin{center}
\scriptsize
\path{72ca9db25cfcd77fdbe95592a2b1b2debbf9bbf006a16a419dd97a31c7b57309}.
\end{center}

Its transaction identifier is:

\begin{center}
\scriptsize
\path{960e5630c0427ca7f74c87930b4611646c9e124969efb31b6514dca5c95be031}.
\end{center}

The transaction was committed in block 6 at \texttt{2026-07-26T13:21:13.565Z}; the channel height was seven after the
run. The event name was \path{M4MetricEvidenceSubmitted}, and the canonical 26-field record SHA-256 was
\path{0f6840a59706211533ac4e84987a26baf89dccce222f8b61ef22bd443ebe3d29}. An independent Python validator
retrieved block 6 through each peer's local QSCC service, found the transaction, and confirmed validation code 0
(\texttt{VALID}). Thus the result is not inferred only from a successful Gateway return.

The independent validation passed all 47 environment, packaging, lifecycle, policy, Gateway-result, dependency, block, and
artifact-integrity checks. As a separate replay negative case, it changed one byte in a temporary copy of the Git bundle. The
copy's SHA-256 changed from
\path{a41205eff782ce4af65f345e0afb97f16fd8677af215c7bdd06f460ce9602272} to
\path{debeb80119a125f617cfdc5ec33537460ab07625a49592711d56601e2f203fdf}, and independent replay rejected it.
The canonical bundle remained unchanged.

\subsection{Interpretation and limits}
\label{sec:phase3_fabric_interpretation}

The network run confirms the implemented mechanism for one pinned evidence fixture: authenticated authorization, exact
arithmetic validation, two-organization endorsement, immutable commit, consistent peer reads, queries, history, event
delivery, and rejection paths work together on the frozen local topology. It does not show that the off-chain numerator and
denominator are semantically true merely because Fabric committed them; that claim remains auditable through the artifact,
policy, verifier, and hashes.

The evaluation uses a single-host educational network, generated test identities, one orderer, one peer per organization,
one fixture, and no load or fault experiment. It therefore supports a bounded proof-of-concept implementation claim, not production
claims about throughput, scale, availability, certificate governance, privacy, or adversarial consensus. Exact commands,
container-image digests, source/output commitments, and machine-readable results are retained in the Phase III
\path{FABRIC_NETWORK_REPORT.md}, \path{fabric_runtime/README.md}, and \path{out/} evidence.
